\chapter{Wstęp, podstawy teoretyczne i metodyka OSINT}\label{ch:wstep}

\section{Cel i zakres projektu}\label{sec:cel-zakres}
\subsection{Uzasadnienie wyboru celu (ETH Zürich - D-INFK)}\label{subsec:wybor-celu}
\subsection{Formalne granice rekonesansu i reguły rozszerzania zakresu}\label{subsec:granice-zakresu}

\section{Podstawy teoretyczne białego wywiadu (OSINT)}\label{sec:podstawy-teoretyczne}
\subsection{Definicja i taksonomia źródeł jawnych w cyberbezpieczeństwie}\label{subsec:definicja-osint}
\subsection{Cykl wywiadowczy w analizie bezpieczeństwa teleinformatycznego}\label{subsec:cykl-wywiadowczy}
\subsection{Różnice operacyjne między rekonesansem pasywnym a aktywnym}\label{subsec:pasywny-vs-aktywny}

\section{Ramy prawne, etyczne i polityka anonimizacji}\label{sec:etyka-anonimizacja}
\subsection{Zgodność z regulaminem projektu i przepisami prawa (FADP / RODO)}\label{subsec:zgodnosc-prawna}
\subsection{Zasady anonimizacji danych osobowych w raporcie i prezentacji}\label{subsec:zasady-anonimizacji}

\section{Warsztat narzędziowy rekonesansu}\label{sec:warsztat-narzedziowy}
\subsection{Frameworki OSINT (osintframework.pl, osintframework.com)}\label{subsec:frameworki-osint}
\subsection{Narzędzia identyfikacji tożsamości cyfrowych (whatsmyname.app)}\label{subsec:whatsmyname}
\subsection{Pasywne narzędzia sieciowe i analityczne}\label{subsec:narzedzia-pasywne}
